\SourceLeaf{12}
itaque punctorum multitudo par erit vel impar, prout $M$ situm est extra soliditatem corporis vel intra, patetque casum ubi $M$ in ipsa corporis superficie iacet, annumerari debere vel casui priori vel posteriori, prout radius $M\Pi$ ab initio vel a corporis soliditate recedit, vel eam intrat. Concipiantur porro rectae a $M$ ad peripheriam spatioli $\mathrm{d}\Sigma$ ductae, quae formabunt superficiem conicam (sensu latiori), atque in superficie corporis nostri ad puncta $P',P'',P'''$ etc. resp. spatiola $\mathrm{d}s',\mathrm{d}s'',\mathrm{d}s'''$ etc. definient. Denique describantur per puncta $P',P'',P'''$ etc. portiunculae superficierum sphaericarum e centro $M$ radiis
\[
MP'=r',\quad MP''=r'',\quad MP'''=r'''\text{ etc.}
\]
sintque spatiola, quae conus ex illis exsecat, $\mathrm{d}\sigma',\mathrm{d}\sigma'',\mathrm{d}\sigma'''$ etc. Omnia haec spatiola $\mathrm{d}\Sigma,\mathrm{d}s',\mathrm{d}\sigma'$ etc. tamquam positiva spectabimus. His praemissis habemus
\[
\mathrm{d}\Sigma=\frac{\mathrm{d}\sigma'}{r'r'}=\frac{\mathrm{d}\sigma''}{r''r''}=\frac{\mathrm{d}\sigma'''}{r'''r'''}\text{ etc.}
\]
Spatiolum $\mathrm{d}\sigma'$ considerari potest tamquam proiectio spatioli $\mathrm{d}s'$ in planum, cui recta $P'M$ est normalis. Hinc erit $\mathrm{d}\sigma'=\pm\mathrm{d}s'.\cos MQ'$, signo superiori vel inferiori accepto, prout $MQ'$ acutus est vel obtusus: casus prior locum habet, quoties recta a $P'$ ad $M$ ducta a corpore recedit, i. e. quoties $M$ iacet extra corpus, casus posterior vero, quoties recta $P'M$ in $P'$ corpus intrat, i. e. quoties $M$ iacet intra corpus. Perinde erit $\mathrm{d}\sigma''=\mp\mathrm{d}s''\cos MQ''$, $\mathrm{d}\sigma'''=\pm\mathrm{d}s'''.\cos MQ'''$ etc. unde patet,
\par
I.\quad Si $M$ iaceat extra corpus, haberi
\[
\begin{aligned}
\mathrm{d}s'.\cos MQ'&=+r'r'\mathrm{d}\Sigma\\
\mathrm{d}s''.\cos MQ''&=-r''r''\mathrm{d}\Sigma\\
\mathrm{d}s'''.\cos MQ'''&=+r'''r'''\mathrm{d}\Sigma\\
&\text{etc.}
\end{aligned}
\]
\par
II.\quad Si vero $M$ iaceat intra corpus, fieri
\[
\begin{aligned}
\mathrm{d}s'.\cos MQ'&=-r'r'\mathrm{d}\Sigma\\
\mathrm{d}s''.\cos MQ''&=+r''r''\mathrm{d}\Sigma\\
\mathrm{d}s'''.\cos MQ'''&=-r'''r'''\mathrm{d}\Sigma\\
&\text{etc.}
\end{aligned}
\]
\par
In casu I itaque erit (propter aequationum multitudinem parem)
